% GENERATED FILE. Edit canonical lessons, then run npm run generate:books.
% canonical-source-hash: fnv1a64:df65fd120147fbcb
% canonical-lessons: GU-C392-sajj, GU-C392-jagrat, GU-C392-upakaran, GU-C392-atmavisvas, GU-C392-gothvavu, GU-R392-review-386-392, GU-R392-review-final-392

\chapter{Ready for Anything}
\label{ch:gu-ready-for-anything}
\hlchaptermodality{392}

\begin{chapteropening}
\textbf{By the end of this chapter:} \emph{I can say how we will be ready if trouble comes.}

Complete the last lesson of chapter 392: I can say how we will be ready if trouble comes.
\end{chapteropening}

\section[\texorpdfstring{\gu{સજ્જ} (sajj)}{સજ્જ (sajj)}]{\gu{સજ્જ} (sajj) — fully equipped; ready}

\label{lesson:GU-C392-sajj}



\begin{quote}
Before the new one: say the Gujarati for to get scattered, then the Gujarati for to get wiped out.
\end{quote}

\subsection*{You'll want to know: \gu{સજ્જ}}
\textbf{\gu{સજ્જ}} — \emph{sajj} — \textquotedblleft{}fully equipped; ready\textquotedblright{}.

\begin{quote}
\gu{જો} \gu{તોફાન} \gu{આવશે} \gu{તો} \gu{પણ} \gu{અમે} \gu{સજ્જ} \gu{છીએ}.
\end{quote}

\begin{quote}
\emph{jo tophān āvśe to paṇ ame sajj chīe.}
\end{quote}

\begin{quote}
Even if a storm comes, we are ready.
\end{quote}

\begin{cousinweb}
\textbf{\gu{સજ્જ}} (\emph{sajj}) is a tatsama, borrowed from Sanskrit \emph{sajja}, \textquotedblleft{}fixed, prepared, equipped\textquotedblright{} (Monier-Williams).
\end{cousinweb}

\begin{grammarlens}[title={What you've built}]
One more word for saying what happens if.
\end{grammarlens}

\subsection*{Your turn}
\begin{itemize}
\raggedright
  \item \emph{Say it:} \emph{sajj}
  \item \emph{Say it:} \emph{sajj}, once more
  \item \emph{Say it:} \emph{jo tophān āvśe to paṇ ame sajj chīe.} — Even if a storm comes, we are ready.
  \item \emph{Recall:} say the Gujarati for to get scattered, then the Gujarati for to get wiped out, then say \emph{sajj} again
\end{itemize}

\begin{tcolorbox}[breakable,colback=teal!4,colframe=teal!35!black,title={Before you move on}]
What does \gu{સજ્જ} mean? (\textquotedblleft{}Fully equipped; ready\textquotedblright{}.) Say it once more.
\end{tcolorbox}
\section[\texorpdfstring{\gu{જાગ્રત} (jāgrat)}{જાગ્રત (jāgrat)}]{\gu{જાગ્રત} (jāgrat) — watchful}

\label{lesson:GU-C392-jagrat}



\begin{quote}
Before the new one: say the Gujarati for to get wiped out, then the Gujarati for fully equipped.
\end{quote}

\subsection*{You'll want to know: \gu{જાગ્રત}}
\textbf{\gu{જાગ્રત}} — \emph{jāgrat} — \textquotedblleft{}watchful\textquotedblright{}.

\begin{quote}
\gu{જો} \gu{ચોકીદાર} \gu{જાગ્રત} \gu{હશે} \gu{તો} \gu{કોઈ} \gu{અંદર} \gu{નહીં} \gu{ઘૂસે}.
\end{quote}

\begin{quote}
\emph{jo cokīdār jāgrat haśe to koī andar nahī̃ ghūse.}
\end{quote}

\begin{quote}
If the watchman is alert, nobody will get in.
\end{quote}

\begin{cousinweb}
\textbf{\gu{જાગ્રત}} (\emph{jāgrat}) is a tatsama, borrowed from Sanskrit \emph{jāgrat}, \textquotedblleft{}waking, watchful\textquotedblright{} (Monier-Williams).
\end{cousinweb}

\begin{grammarlens}[title={What you've built}]
One more word for saying what happens if.
\end{grammarlens}

\subsection*{Your turn}
\begin{itemize}
\raggedright
  \item \emph{Say it:} \emph{jāgrat}
  \item \emph{Say it:} \emph{jāgrat}, once more
  \item \emph{Say it:} \emph{jo cokīdār jāgrat haśe to koī andar nahī̃ ghūse.} — If the watchman is alert, nobody will get in.
  \item \emph{Recall:} say the Gujarati for to get wiped out, then the Gujarati for fully equipped, then say \emph{jāgrat} again
\end{itemize}

\begin{tcolorbox}[breakable,colback=teal!4,colframe=teal!35!black,title={Before you move on}]
What does \gu{જાગ્રત} mean? (\textquotedblleft{}Watchful\textquotedblright{}.) Say it once more.
\end{tcolorbox}
\section[\texorpdfstring{\gu{ઉપકરણ} (upakaraṇ)}{ઉપકરણ (upakaraṇ)}]{\gu{ઉપકરણ} (upakaraṇ) — a device}

\label{lesson:GU-C392-upakaran}



\begin{quote}
Before the new one: say the Gujarati for fully equipped, then the Gujarati for watchful.
\end{quote}

\subsection*{You'll want to know: \gu{ઉપકરણ}}
\textbf{\gu{ઉપકરણ}} — \emph{upakaraṇ} — \textquotedblleft{}a device\textquotedblright{}.

\begin{quote}
\gu{જો} \gu{વીજળી} \gu{જશે} \gu{તો} \gu{આ} \gu{ઉપકરણ} \gu{કામ} \gu{નહીં} \gu{કરે}.
\end{quote}

\begin{quote}
\emph{jo vījḷī jaśe to ā upakaraṇ kām nahī̃ kare.}
\end{quote}

\begin{quote}
If the power goes, this device will not work.
\end{quote}

A neuter noun.

\begin{cousinweb}
\textbf{\gu{ઉપકરણ}} (\emph{upakaraṇ}) is a tatsama, borrowed from Sanskrit \emph{upakaraṇa}, \textquotedblleft{}doing a service; an implement, a means\textquotedblright{} (Monier-Williams).
\end{cousinweb}

\begin{grammarlens}[title={What you've built}]
One more word for saying what happens if.
\end{grammarlens}

\subsection*{Your turn}
\begin{itemize}
\raggedright
  \item \emph{Say it:} \emph{upakaraṇ}
  \item \emph{Say it:} \emph{upakaraṇ}, once more
  \item \emph{Say it:} \emph{jo vījḷī jaśe to ā upakaraṇ kām nahī̃ kare.} — If the power goes, this device will not work.
  \item \emph{Recall:} say the Gujarati for fully equipped, then the Gujarati for watchful, then say \emph{upakaraṇ} again
\end{itemize}

\begin{tcolorbox}[breakable,colback=teal!4,colframe=teal!35!black,title={Before you move on}]
What does \gu{ઉપકરણ} mean? (\textquotedblleft{}A device\textquotedblright{}.) Say it once more.
\end{tcolorbox}
\section[\texorpdfstring{\gu{આત્મવિશ્વાસ} (ātmaviśvās)}{આત્મવિશ્વાસ (ātmaviśvās)}]{\gu{આત્મવિશ્વાસ} (ātmaviśvās) — self-confidence}

\label{lesson:GU-C392-atmavisvas}



\begin{quote}
Before the new one: say the Gujarati for watchful, then the Gujarati for a device.
\end{quote}

\subsection*{You'll want to know: \gu{આત્મવિશ્વાસ}}
\textbf{\gu{આત્મવિશ્વાસ}} — \emph{ātmaviśvās} — \textquotedblleft{}self-confidence\textquotedblright{}.

\begin{quote}
\gu{જો} \gu{તમારામાં} \gu{આત્મવિશ્વાસ} \gu{હશે} \gu{તો} \gu{તમે} \gu{કોઈ} \gu{પણ} \gu{મુશ્કેલીને} \gu{પહોંચી} \gu{વળશો}.
\end{quote}

\begin{quote}
\emph{jo tamārāmā̃ ātmaviśvās haśe to tame koī paṇ muśkelīne pahõcī vaḷśo.}
\end{quote}

\begin{quote}
If you believe in yourself, you will cope with any difficulty.
\end{quote}

A masculine noun.

\begin{cousinweb}
\textbf{\gu{આત્મવિશ્વાસ}} (\emph{ātmaviśvās}) is built from Sanskrit \emph{ātman}, \textquotedblleft{}self\textquotedblright{}, and \emph{viśvāsa}, \textquotedblleft{}trust\textquotedblright{} (Monier-Williams): trust in oneself. \textbf{\gu{વિશ્વાસ}} (\emph{viśvās}) is the same second word.
\end{cousinweb}

\begin{grammarlens}[title={What you've built}]
One more word for saying what happens if.
\end{grammarlens}

\subsection*{Your turn}
\begin{itemize}
\raggedright
  \item \emph{Say it:} \emph{ātmaviśvās}
  \item \emph{Say it:} \emph{ātmaviśvās}, once more
  \item \emph{Say it:} \emph{jo tamārāmā̃ ātmaviśvās haśe to tame koī paṇ muśkelīne pahõcī vaḷśo.} — If you believe in yourself, you will cope with any difficulty.
  \item \emph{Recall:} say the Gujarati for watchful, then the Gujarati for a device, then say \emph{ātmaviśvās} again
\end{itemize}

\begin{tcolorbox}[breakable,colback=teal!4,colframe=teal!35!black,title={Before you move on}]
What does \gu{આત્મવિશ્વાસ} mean? (\textquotedblleft{}Self-confidence\textquotedblright{}.) Say it once more.
\end{tcolorbox}
\section[\texorpdfstring{\gu{ગોઠવાવું} (goṭhvāvũ)}{ગોઠવાવું (goṭhvāvũ)}]{\gu{ગોઠવાવું} (goṭhvāvũ) — to fall into place}

\label{lesson:GU-C392-gothvavu}



\begin{quote}
Before the new one: say the Gujarati for a device, then the Gujarati for self-confidence.
\end{quote}

\subsection*{You'll want to know: \gu{ગોઠવાવું}}
\textbf{\gu{ગોઠવાવું}} — \emph{goṭhvāvũ} — \textquotedblleft{}to fall into place\textquotedblright{}.

\begin{quote}
\gu{જો} \gu{બધા} \gu{મદદ} \gu{કરશે} \gu{તો} \gu{બધું} \gu{જલદી} \gu{ગોઠવાઈ} \gu{જશે}.
\end{quote}

\begin{quote}
\emph{jo badhā madad karśe to badhũ jaldī goṭhvāī jaśe.}
\end{quote}

\begin{quote}
If everyone helps, everything will soon fall into place.
\end{quote}

Past (he …): \gu{તે} \gu{નવા} \gu{ઘરમાં} \gu{ગોઠવાઈ} \gu{ગયો}, \emph{te navā gharmā̃ goṭhvāī gayo}.

\begin{cousinweb}
\textbf{\gu{ગોઠવાવું}} (\emph{goṭhvāvũ}) is the \emph{-āvũ} partner of \textbf{\gu{ગોઠવવું}} (\emph{goṭhavvũ}), \textquotedblleft{}to arrange\textquotedblright{}, whose own history is not settled.
\end{cousinweb}

\begin{grammarlens}[title={What you've built}]
One more word for saying what happens if.
\end{grammarlens}

\subsection*{Your turn}
\begin{itemize}
\raggedright
  \item \emph{Say it:} \emph{goṭhvāvũ}
  \item \emph{Say it:} \emph{goṭhvāvũ}, once more
  \item \emph{Say it:} \emph{jo badhā madad karśe to badhũ jaldī goṭhvāī jaśe.} — If everyone helps, everything will soon fall into place.
  \item \emph{Recall:} say the Gujarati for a device, then the Gujarati for self-confidence, then say \emph{goṭhvāvũ} again
\end{itemize}

\begin{tcolorbox}[breakable,colback=teal!4,colframe=teal!35!black,title={Before you move on}]
What does \gu{ગોઠવાવું} mean? (\textquotedblleft{}To fall into place\textquotedblright{}.) Say it once more.
\end{tcolorbox}
\section[(dialogue)]{Looking back: from Trust and Dependence to Ready for Anything}

\label{lesson:GU-R392-review-386-392}



\begin{quote}
Say the Gujarati for belief and a refuge.
\end{quote}

\subsection*{Your turn}
Say these aloud:

\begin{itemize}
\raggedright
  \item belief, a refuge, insistence, subject to, to give in
  \item better, excellent, superior, lovely, to deteriorate
  \item an outline, an initiative, a public drive, to commence, to undertake
  \item a decision, one's own will, a dilemma, to hesitate, to waver
  \item a crop, mature, barren, to grow plants, to be produced
  \item a mix-up, disorder, to shatter, to get scattered, to get wiped out
  \item fully equipped, watchful, a device, self-confidence, to fall into place
\end{itemize}

\begin{tcolorbox}[breakable,colback=teal!4,colframe=teal!35!black,title={Before you move on}]
Belief? (\textbf{\gu{આસ્થા}}, \emph{āsthā}.) A refuge? (\textbf{\gu{આશરો}}, \emph{āśro}.) Insistence? (\textbf{\gu{આગ્રહ}}, \emph{āgrah}.) Subject to? (\textbf{\gu{આધીન}}, \emph{ādhīn}.) To give in? (\textbf{\gu{ઝૂકવું}}, \emph{jhūkvũ}.) Better? (\textbf{\gu{બહેતર}}, \emph{bahetar}.) Excellent? (\textbf{\gu{ઉત્તમ}}, \emph{uttam}.) Superior? (\textbf{\gu{ચડિયાતું}}, \emph{caḍiyātũ}.) Lovely? (\textbf{\gu{સરસ}}, \emph{saras}.) To deteriorate? (\textbf{\gu{કથળવું}}, \emph{kathaḷvũ}.) An outline? (\textbf{\gu{રૂપરેખા}}, \emph{rūprekhā}.) An initiative? (\textbf{\gu{ઉપક્રમ}}, \emph{upakram}.) A public drive? (\textbf{\gu{અભિયાન}}, \emph{abhiyān}.) To commence? (\textbf{\gu{આરંભવું}}, \emph{ārambhvũ}.) To undertake? (\textbf{\gu{આદરવું}}, \emph{ādarvũ}.) A decision? (\textbf{\gu{નિર્ણય}}, \emph{nirṇay}.) One's own will? (\textbf{\gu{મરજી}}, \emph{marjī}.) A dilemma? (\textbf{\gu{દ્વિધા}}, \emph{dvidhā}.) To hesitate? (\textbf{\gu{ખચકાવું}}, \emph{khackāvũ}.) To waver? (\textbf{\gu{ડગવું}}, \emph{ḍagvũ}.) A crop? (\textbf{\gu{પાક}}, \emph{pāk}.) Mature? (\textbf{\gu{પરિપક્વ}}, \emph{paripakva}.) Barren? (\textbf{\gu{ઉજ્જડ}}, \emph{ujjaḍ}.) To grow plants? (\textbf{\gu{ઉગાડવું}}, \emph{ugāḍvũ}.) To be produced? (\textbf{\gu{ઊપજવું}}, \emph{ūpajvũ}.) A mix-up? (\textbf{\gu{ગડબડ}}, \emph{gaḍbaḍ}.) Disorder? (\textbf{\gu{અવ્યવસ્થા}}, \emph{avyavasthā}.) To shatter? (\textbf{\gu{ભાંગવું}}, \emph{bhā̃gvũ}.) To get scattered? (\textbf{\gu{વેરાવું}}, \emph{verāvũ}.) To get wiped out? (\textbf{\gu{ભૂંસાવું}}, \emph{bhū̃sāvũ}.) Fully equipped? (\textbf{\gu{સજ્જ}}, \emph{sajj}.) Watchful? (\textbf{\gu{જાગ્રત}}, \emph{jāgrat}.) A device? (\textbf{\gu{ઉપકરણ}}, \emph{upakaraṇ}.) Self-confidence? (\textbf{\gu{આત્મવિશ્વાસ}}, \emph{ātmaviśvās}.) To fall into place? (\textbf{\gu{ગોઠવાવું}}, \emph{goṭhvāvũ}.)
\end{tcolorbox}
\section[(dialogue)]{Ready for Anything — once more}

\label{lesson:GU-R392-review-final-392}



\begin{quote}
Say the Gujarati for fully equipped and watchful.
\end{quote}

\subsection*{Your turn}
\begin{itemize}
\raggedright
  \item \emph{Say it:} fully equipped, watchful, a device, self-confidence, to fall into place
\end{itemize}

\begin{tcolorbox}[breakable,colback=teal!4,colframe=teal!35!black,title={Before you move on}]
Fully equipped? (\textbf{\gu{સજ્જ}}, \emph{sajj}.) Watchful? (\textbf{\gu{જાગ્રત}}, \emph{jāgrat}.) A device? (\textbf{\gu{ઉપકરણ}}, \emph{upakaraṇ}.) Self-confidence? (\textbf{\gu{આત્મવિશ્વાસ}}, \emph{ātmaviśvās}.) To fall into place? (\textbf{\gu{ગોઠવાવું}}, \emph{goṭhvāvũ}.)
\end{tcolorbox}
